\documentclass[11pt,a4paper]{article}

\usepackage[T1]{fontenc}
\usepackage[utf8]{inputenc}
\usepackage{lmodern}
\usepackage[english]{babel}
\usepackage[
  top=2.8cm,
  bottom=2.8cm,
  left=3cm,
  right=3cm,
  headheight=14pt
]{geometry}
\usepackage{microtype}
\usepackage{graphicx}
\usepackage{booktabs}
\usepackage{tabularx}
\usepackage{float}
\newcolumntype{Y}{>{\raggedright\arraybackslash}X}
\usepackage{enumitem}
\usepackage{hyperref}
\usepackage{xcolor}
\usepackage{listings}
\usepackage{fancyhdr}
\usepackage{titlesec}
\usepackage{tcolorbox}
\tcbuselibrary{skins,breakable}

\definecolor{brand}{RGB}{0,82,147}
\definecolor{muted}{RGB}{96,96,96}

\hypersetup{
  colorlinks=true,
  linkcolor=brand,
  citecolor=brand,
  urlcolor=brand,
  pdftitle={libp2p-raft Whitepaper},
}

\pagestyle{fancy}
\fancyhf{}
\fancyhead[L]{\small\textcolor{muted}{libp2p-raft Whitepaper}}
\fancyhead[R]{\small\textcolor{muted}{\thepage}}
\renewcommand{\headrulewidth}{0.4pt}

\titleformat{\section}{\large\bfseries\color{brand}}{\thesection}{1em}{}
\titleformat{\subsection}{\normalsize\bfseries}{\thesubsection}{1em}{}

\lstset{
  basicstyle=\ttfamily\footnotesize,
  breaklines=true,
  columns=fullflexible,
  frame=single,
  rulecolor=\color{black!15},
  backgroundcolor=\color{black!3},
}

\newtcolorbox{designbox}[1][]{
  enhanced,
  breakable,
  colback=brand!4,
  colframe=brand!60,
  fonttitle=\bfseries,
  title=#1,
  boxrule=0.6pt,
  arc=2pt,
}

% --- Title block (whitepaper style) ---
\newcommand{\WhitepaperTitle}{%
  \begin{titlepage}
    \centering
    \vspace*{1.5cm}
    {\LARGE\bfseries\color{brand} WHITEPAPER\par}
    \vspace{0.6cm}
    {\Huge\bfseries A minimal Raft consensus layer\\on rust-libp2p\par}
    \vspace{0.4cm}
    {\large\textit{libp2p-raft --- Swarm, Behaviour, Handler architecture}\par}
    \vspace{2cm}
    {\large HoodLabs / libp2p-raft\par}
    \vspace{0.3cm}
    {\normalsize\textcolor{muted}{Document version: 1.0 \quad|\quad \today}\par}
    \vfill
    {\small\textcolor{muted}%
      This document describes a research and experimentation design.\\
      It is not intended for production deployment.\par}
  \end{titlepage}
}

% Block paragraphs (no first-line indent) and avoid overfull lines from \texttt
\setlength{\parindent}{0pt}
\setlength{\parskip}{0.5em plus 0.2em minus 0.1em}
\setlength{\emergencystretch}{3em}
\frenchspacing
% Lists are code-heavy: ragged-right avoids stretched gaps and overflow
\setlist{before=\raggedright}

% Full-width abstract, aligned with body text
\renewenvironment{abstract}{%
  \begin{center}\bfseries\abstractname\end{center}\vspace{-0.5em}\noindent\ignorespaces
}{\par\medskip}

% Inline code: allow a line break after each underscore
\newcommand{\code}[1]{{\def\_{\textunderscore\allowbreak}\texttt{#1}}}

\begin{document}
\WhitepaperTitle

\begin{abstract}
\noindent
We present \emph{libp2p-raft}: a research implementation that embeds a simple Raft state machine
into a rust-libp2p \code{NetworkBehaviour}. The design separates \textbf{RaftEngine}
(pure logic, synchronous, no I/O) from \textbf{RaftBehaviour} (the network adapter: maps
\code{NodeId} $\leftrightarrow$ \code{PeerId}, manages asynchronous RPCs) and
\textbf{RaftHandler} (one instance per connection, unary framed RPC over the
\code{/libp2p-raft/1.0.0} substream). The libp2p Swarm orchestrates transports, the connection
pool and the event loop; behaviour and handler never call each other directly.

The paper states the architectural goals, the trust model, the wire format, the control flows
(leader election, replication, inbound reply), the event-driven execution model, and the
deliberate limitations compared with production Raft. It is intended as a reference for learning
and experimenting with combining consensus and a modern P2P stack.
\end{abstract}

\noindent\textbf{Keywords:} Raft, libp2p, rust-libp2p, Swarm, NetworkBehaviour,
ConnectionHandler, distributed consensus, unary RPC.

\tableofcontents
\newpage

\section{Introduction}

A replicated state machine needs a consensus protocol so that processes agree on an ordered
history of commands despite network failures or crashes. Raft~\cite{ongaro2014raft} splits the
problem into leader election, log replication and safety invariants that are easier to implement
than Paxos. In parallel, libp2p~\cite{libp2p} provides a modular networking stack: transport
upgrades, cryptographic identity, substream multiplexing --- applications express protocols via
\code{NetworkBehaviour} and \code{ConnectionHandler} instead of raw sockets.

\code{libp2p-raft} explores a two-layer \emph{composition}: a hand-written Raft engine owned by
\code{RaftBehaviour}, which neither dials nor serializes by itself; byte framing lives in
\code{RaftHandler}. The library does not own the \code{Swarm}: the host application assembles
TCP, Noise and Yamux via \code{SwarmBuilder} and attaches \code{RaftBehaviour}.

\section{Problem and design goals}

\subsection{Problem}

Given a fixed set of voters (static configuration), the nodes must:
\begin{enumerate}[label=\roman*., nosep]
  \item elect at most one leader per term;
  \item replicate the command log from the leader;
  \item commit and apply the same log prefix on a majority;
  \item communicate over an authenticated (Noise) and multiplexed (Yamux) P2P network.
\end{enumerate}

\subsection{Architectural goals}

\begin{designbox}[Design principles]
\begin{itemize}[nosep]
  \item \textbf{Separate engine and network:} the engine only receives ticks, logical RPCs and
        proposals, and returns \code{Action}s --- no \code{PeerId}, no stream handles.
  \item \textbf{Unary RPC over substreams:} each RPC opens a substream, exchanges one
        request/response, and closes it --- avoiding head-of-line blocking between unrelated RPCs.
  \item \textbf{The Swarm is the intermediary:} behaviour $\leftrightarrow$ handler via
        \code{FromBehaviour} / \code{ToBehaviour}; no direct calls.
  \item \textbf{Lost connection $\neq$ lost membership:} quorum follows the voter configuration;
        redials and RPC timeouts never remove a voter.
\end{itemize}
\end{designbox}

\subsection{Out of scope}

The implementation does \textbf{not} target production: storage may be in-memory or a simple
file; AppendEntries pipelining is limited to depth 1; joint consensus and full snapshot support are
incomplete. The value lies in the \emph{layering} and the Behaviour--Handler lifecycle.

\section{System overview}

\subsection{Layers}

\begin{table}[H]
\centering
\caption{Layer responsibilities}
\begin{tabularx}{\textwidth}{@{}lY@{}}
\toprule
\textbf{Layer} & \textbf{Responsibility} \\
\midrule
Application / Swarm & Build transports, \code{listen\_on}, dial seeds, \code{propose}, receive
\code{SwarmEvent::Behaviour} \\
\code{RaftBehaviour} & Engine, \code{PeerMap}, \code{pending}, \code{inbound\_channels},
timer, \code{execute\_actions} \\
\code{RaftHandler} & Substreams, negotiate \code{/libp2p-raft/1.0.0}, encode/decode frames \\
\code{RaftEngine} & \code{tick}, \code{handle\_rpc}, election, replication, commit \\
\code{Storage} & Hard state, log, (snapshot) \\
\bottomrule
\end{tabularx}
\end{table}

\subsection{The Swarm in libp2p}

The \textbf{Swarm} is the node-level event loop: it owns transports, listeners, dialing and the
connection pool (one \code{ConnectionHandler} per connection). The name \emph{swarm} evokes a P2P
swarm of peers --- a node managing a dynamic set of connections, not one fixed client--server channel.

The application drives \code{swarm.select\_next\_some()}. The Swarm polls the behaviour and the
handlers; Raft logic does not run inside the Swarm.

\subsection{Logical flow diagram}

\noindent\fbox{\parbox{\dimexpr\textwidth-2\fboxsep-2\fboxrule\relax}{\small\ttfamily
App $\rightarrow$ propose / events\\
\quad $\downarrow$\\
Swarm $\leftrightarrow$ RaftBehaviour $\leftrightarrow$ RaftEngine $\leftrightarrow$ Storage\\
\quad $\downarrow$ NotifyHandler / HandlerEvent\\
RaftHandler $\leftrightarrow$ Yamux substream $\leftrightarrow$ peer
}}

\section{Behaviour and Handler interfaces}

\subsection{Associated types}

The \code{ConnectionHandler} trait requires two message types (the project defines their content):

\begin{itemize}[nosep]
  \item \code{FromBehaviour}: behaviour $\rightarrow$ handler --- \code{SendRequest},
        \code{SendResponse}.
  \item \code{ToBehaviour}: handler $\rightarrow$ behaviour --- \code{Request},
        \code{Response}, \code{Failure}.
\end{itemize}

\subsection{Outbound and inbound unary}

\textbf{Outbound:}
\begin{enumerate}[label=\arabic*., nosep]
  \item The behaviour calls \code{send\_rpc} $\rightarrow$ \code{SendRequest}.
  \item The handler opens an \emph{outbound substream} on the connection to the peer.
  \item It writes the frame and reads the response $\rightarrow$ \code{ToBehaviour::Response}.
\end{enumerate}

\textbf{Inbound:}
\begin{enumerate}[label=\arabic*., nosep]
  \item The peer opens a substream; the handler reads the frame.
  \item \code{ToBehaviour::Request} with a \code{channel\_id} $\rightarrow$ the engine processes it.
  \item The reply goes back via \code{SendResponse} on the \emph{same} substream.
\end{enumerate}

\subsection{Replying on the same connection: \code{to == from}}

When \code{execute\_actions} is called with \code{Some((from, channel\_id))}:
\begin{itemize}[nosep]
  \item If \code{Action::Send \{ to, msg \}} has \code{to == from}, the behaviour uses
    \code{NotifyHandler::One(connection\_id)} and \code{SendResponse} --- no new outbound RPC is opened.
  \item \code{inbound\_channels} stores \code{channel\_id} $\rightarrow$ \code{\{peer, conn, correlation\_id\}}.
\end{itemize}

\code{NotifyHandler::Any} is only for new outbound RPCs: the Swarm picks \textbf{one} connection
to the peer. \code{pending} maps the \code{correlation\_id} of RPCs \textbf{we sent} that are
awaiting a response.

\section{Wire protocol specification}

\subsection{Substream identifier}

Multistream-select negotiates the protocol:
\begin{lstlisting}
/libp2p-raft/1.0.0
\end{lstlisting}

\subsection{Framing}

Each message on a substream:
\begin{enumerate}[nosep]
  \item a 4-byte payload length, big-endian (\code{u32});
  \item the payload \code{bincode(WireEnvelope)} with
        \code{WireEnvelope \{ correlation\_id, msg \}};
  \item \code{msg}: a \code{RaftMessage} variant (RequestVote, AppendEntries, snapshot, ...).
\end{enumerate}

The handler enforces a frame size limit (e.g.\ 4\,MiB). A Raft heartbeat is an
\code{AppendEntries} with empty \code{entries}.

\subsection{In-memory decoding}

\code{decode\_envelope} wraps the \code{\&[u8]} in a \code{std::io::Cursor} to read the length
prefix through the \code{Read} API, then deserializes the rest --- no file involved.

\section{Consensus control flows}

\subsection{Leader election}

Follower/Candidate: \code{election\_timeout} expires $\rightarrow$ \code{engine.tick} $\rightarrow$
increment term, vote for self, \code{Broadcast RequestVote} $\rightarrow$ the behaviour opens an RPC to
each voter. Majority $\rightarrow$ \code{BecomeLeader}, initialise \code{next\_index} /
\code{match\_index}, start heartbeats.

\subsection{Heartbeat and replication}

Only the \textbf{Leader} sends \code{AppendEntries} every \code{heartbeat\_interval} via
\code{replicate\_actions}. The engine emits actions $\rightarrow$ behaviour \code{send\_rpc}
$\rightarrow$ handler. Followers reset their election timeout on each heartbeat.

\subsection{Proposing and committing commands}

Leader: \code{propose(data)} persists the entry $\rightarrow$ replication $\rightarrow$ majority
\code{match\_index} $\rightarrow$ \code{commit\_index} $\rightarrow$ \code{Apply} $\rightarrow$
a \code{Committed} event for the application.

\subsection{Timeouts and transport errors}

An outbound RPC past \code{rpc\_timeout}: fail the pending entry, the engine clears
\code{ae\_inflight}; retries happen implicitly through tick/election/heartbeat. Dial failure:
backoff, never treated as removing a voter.

\section{Execution model}

A single Swarm event loop (tokio). The engine is called \textbf{synchronously} from the behaviour ---
there is no separate consensus thread. The handler polls many substream I/O operations concurrently
(\code{FuturesUnordered}); the engine remains sequential. Responses may arrive out of send order;
they are matched by \code{correlation\_id}.

The behaviour arms a \code{Sleep} until the earliest deadline (engine or pending RPC) and only
calls \code{tick} when it is due.

\section{Identity and bootstrap}

\begin{itemize}[nosep]
  \item \code{NodeId} (\code{u64}): Raft membership, fixed in the cluster config.
  \item \code{PeerId}: libp2p identity from a keypair; the seed list maps \code{NodeId} $\leftrightarrow$
        \code{PeerId} + multiaddr.
  \item Voters must be configured before the first election timeout. No DHT/gossip is required for a
        minimal deployment.
\end{itemize}

\section{Limitations and implementation status}

\begin{table}[H]
\centering
\caption{Compared with production Raft}
\begin{tabularx}{\textwidth}{@{}lY@{}}
\toprule
\textbf{Implemented} & Election, replication, propose, commit events, multi-node example \\
\textbf{Limitations} & AE pipeline depth 1; reject $\rightarrow$ simple \code{next\_index} decrement;
durability depends on \code{Storage}; full snapshot/membership support incomplete \\
\bottomrule
\end{tabularx}
\end{table}

\section{Conclusion}

\code{libp2p-raft} illustrates a layered architecture: Raft invariants in a pure engine, rust-libp2p
providing authenticated and multiplexed connections, and an adapter that uses correlation IDs, a peer
map and deadline-driven ticks, and distinguishes loss of connectivity from membership change. Within
the simplified scope described, the design is a compact laboratory for studying how to map a consensus
state machine onto a modern P2P stack.

\begin{thebibliography}{9}
\bibitem{ongaro2014raft}
D.~Ongaro and J.~Ousterhout,
\newblock In Search of an Understandable Consensus Algorithm,
\newblock in \emph{Proc.\ USENIX ATC}, 2014.

\bibitem{libp2p}
Protocol Labs and contributors,
\newblock \emph{rust-libp2p} --- Swarm, \code{NetworkBehaviour},
\code{ConnectionHandler}.
\end{thebibliography}

\end{document}
